\documentclass[10pt,a4paper]{amsart}
\usepackage[margin=19mm]{geometry}
\usepackage{amsmath,amssymb,mathtools}
\usepackage[scr=rsfs,cal=euler]{mathalfa}
\usepackage{xfrac}
\usepackage[unicode,hidelinks,bookmarks=false]{hyperref}
\newcommand{\ie}{i.e.\ }
\newcommand{\cf}{cf.\ }
\newcommand{\resp}{respectively\ }
\newcommand{\Subsection}{\subsection}
\providecommand{\nobreakdash}{\nobreak\hbox{-}}
\setlength{\emergencystretch}{2em}
\multlinegap0pt
\usepackage{bm}
\usepackage{bbm}
\usepackage{dsfont}
\usepackage{xspace}
\usepackage{cancel}
\usepackage{mathtools}
\usepackage{amsthm}
\makeatletter
\@ifundefined{theorem}{\newtheorem{theorem}{Theorem}}{}
\makeatother
\newcommand{\lcm}{\mathrm{lcm}}
\title[On Some Properties of LCM-Lattices of Edge Ideals of k-Unifo]{On Some Properties of LCM-Lattices of Edge Ideals of k-Uniform Hypergraphs}
\author{Muneeba Mansha, Sarfraz Ahmad}
\date{Source: arXiv:2605.13617v2 (CC BY 4.0)}
\begin{document}
\maketitle

\noindent\emph{Source and licence.} On Some Properties of LCM-Lattices of Edge Ideals of k-Uniform Hypergraphs. Version arXiv:2605.13617v2 (CC BY 4.0), distributed under the Creative Commons Attribution 4.0 International licence, as recorded in the arXiv licence metadata for this exact version and shown on the abstract page https://arxiv.org/abs/2605.13617v2. Full rights notice: licenses/arxiv-2605.13617-CC-BY-4.0.txt.\par\smallskip

\noindent\textit{Editorial scope.} Counted unit: the Theorem whose source label is \texttt{None} in the article (source0.tex lines 822--824, source label \texttt{None}), delivered together with the article's own complete original proof of it (source0.tex lines 825--854, one \texttt{proof} environment of 30 lines). No mathematics is written, repaired or re-derived: statement and proof are the article's own text line for line. Required context retained uncounted: none. Every definition, display and label of those lines is retained unchanged. Omitted from this package and not redistributed: 1--821, 855--2065. Nothing inside a retained excerpt is omitted or altered: every retained body is delivered exactly as the article's lines are written, so no omission receipt of any kind accompanies this package. Citations: the retained excerpts carry no citation command at all, so no bibliography entry of the article is transcribed and no reference is redistributed. Reference adaptation: none was needed and none was made; every reference inside the delivered unit resolves inside it, so no unresolved reference or citation is produced. Printed slips kept literal: no printed word, symbol, hypothesis or number of the counted statement or of its proof was altered. Layout and class: the article's own class is \texttt{amsart} with packages including amssymb, amsthm, amsfonts, latexsym, mathtools, amsmath, mathrsfs, stmaryrd, inputenc, color, enumitem, accents, hyperref, cleveref; the wrapper is the lane's pinned \texttt{amsart} editorial wrapper at 10pt on A4, which restates the theorem-like environments the retained text needs and adds the article's own macro definitions that the retained text needs (\texttt{\textbackslash lcm}), with extra wrapper packages (bm, bbm, dsfont, xspace, cancel, mathtools). The delivered numbering is the unit's own. Assets: no retained span contains a figure environment, a graphics-inclusion command, a PDF-page inclusion, a table or a TikZ picture; no image file is copied into this package, the package declares no asset dependency and no image file is redistributed with it. Licence: version arXiv:2605.13617v2 of this article is distributed under the Creative Commons Attribution 4.0 International licence, as recorded in the arXiv licence metadata for this exact version and shown on the abstract page https://arxiv.org/abs/2605.13617v2. A full rights notice is delivered at licenses/arxiv-2605.13617-CC-BY-4.0.txt. Characters: every retained excerpt is pure ASCII, so the delivered engine, XeLaTeX with its default Latin Modern text font, reports no missing character.
\begin{theorem}
     Let $G$ be a connected graph and $I(G)$ be its edge ideal, then $lcm(I(G))$ is \emph{relatively complemented} if and only if there exists no induced subgraph $G|_W$ on the vertex set $W \subseteq V$ such that ${G|_W}$ is of path length $ \geq 3$.
\end{theorem}

\begin{proof}
     %  Contrary suppose that there exists a path of length greater than equal 3. Suppose that we have a path $\{x_1,x_2\},\{x_2,x_3\},\{x_3,x_4\},\ldots,\{x_{i-1},x_i\}$ such that $d_i=1$ and $d_{i-1}=2$ then $x_1,x_2,\ldots, x_{i-2}$ will not have a complement because its complement must contains $x_{i-1}$ and $x_i$ that do not exist in the $\lcm$ lattice of this path graph.  So there exists a chain having $x_1,x_2,\ldots, x_{i-2} \longrightarrow x_1,x_2,\ldots, x_{i-1} \longrightarrow x_1,x_2,\ldots, x_{i}$ but here we can not find an element/ atom $x_{i-1}x_i$. So our assumption for being length greater than equal 3 to be complemented is wrong. which implies that a graph with length less than equal 3 is always relatively complemented.\\
%
     %   \textbf{Conversally,} suppose that for a connected graph G and its edge ideal $I(G)$, $lcm(I(G))$ is relatively complemented. We need to prove that there exists no induced subgraph $G| _w$
    %   on the vertex set $W \subseteq V$ such that ${G| _w}$ is a path of length $i \geq 3$.\\
%
    %    Let $u \in \mathcal{L}=\lcm (I(G))$. Let $\supp(u) =A\subseteq V(G)$, such that $A=\cup_{e\in E' \subset E(G)} e.$ Set $B =V(G) \setminus A $ and $A \subset B $  such that chain $x_{i-1}$ %\subset
    %    is properly contained in the chain of elements $x_i$, for all $i$. Then we have two cases:\\
%
%
    % Case $(i):$ If we will choose an interval $[x_1,x_n]$ such that $x_j \in$  $[x_1,x_n]$ such that  then each atom will contain $x_1,x_n$ a path having a shape of  
   % Case $(ii):$

%%%%%%%%%%%%%%%%%%%%%%%%%%%%%%%%
%
    %Let $G$ be a connected graph and $I(G)$ be its edge ideal then 
     Suppose that $\lcm(I(G))$ is relatively complemented. 
To the contrary, assume that there exists an induced subgraph $W \subseteq V$ having a path of length $3$, where 
$W = \{x_1x_2,\, x_2x_3,\, x_3x_4\}$. 
Then, the interval 
\[
[a,\, \{x_1x_2x_3x_4\}]
\]
is not complemented for $a \in \{x_1x_2,\, x_3x_4\}$. \\[4pt]
Conversely, suppose that there exists no induced subgraph of path length $3$. 
This means that $G$ has a maximum induced path of length two. 
We now show that $\lcm(I(G))$ is relatively complemented. 
Since the $\lcm$-lattice of a graph having path length at most $2$ is Boolean, and every Boolean lattice is relatively complemented, we conclude that $\lcm(I(G))$ is relatively complemented.

\end{proof}

\end{document}
